% Opening of the old chapter c25, its \chapter command removed; generated
\section*{Extreme values, hitting times and a path to RH}\label{chap:evt}\label{ch:evt}

The ferromagnetic approximants above are laws whose transforms are proved to have only real zeros; the extreme value approximants below are explicit kernels whose real zeros are tracked to a stated height. The P\'olya density is an exact superposition of extreme value laws (Section~\ref{ev:sec:decomp}), and the reciprocal $\xi$ function is the Laplace transform of a sum of independent exponentials with rates $\gamma_k^2$, which Kent's theorem reads as the hitting time of a one dimensional diffusion; by Krein's inverse spectral theorem every strictly increasing sequence of positive rates arises this way. Each extreme value approximation therefore defines a candidate hitting time, and real simple zeros become the existence of a string with positive mass. None of this proves RH. It gives a precise target, three routes to it with the place where each stops, the folded Gumbel kernels, whose zeros are tracked rigorously to a stated height, and diagnostics, from Hankel determinants of the kernel moments to threshold fits of the hitting time tail, whose power and limits are measured.

This part is based on \cite{P26evt}; its numbers were recomputed, and values not recomputed are marked unchecked. The kernel $\Phi$ here is $\Phi_{\mathrm{dBN}}$ of Section~\ref{ch:dbn} (the normalisation of Rodgers and Tao, with $e^{9u}$ and $e^{4u}$), not the kernel of \eqref{eq:ch2:Phi}; with it $H(z)=\int_0^\infty\Phi(u)\cos(zu)\,du=\tfrac18\Xi(z/2)$, so the zeros of $H$ are at $z=\pm2\gamma_k$ and the variable of the zeros is $x=z/2$. The eigenvalues of a string are $\lambda_k$, the flow time is $t$ and $\Lambda$ the de Bruijn Newman constant (Section~\ref{ch:dbn}), and Euler's constant is $\gamma_E$.

Results proved elsewhere are cited: Kent's theorem and the reduction of RH to a string with simple spectrum (Propositions~\ref{ht:prop:kent} and~\ref{ht:prop:red}, Section~\ref{sec:ch18:kent}), the B spline form of the truncated hitting time (Section~\ref{ht:sec:ev}, Proposition~\ref{ht:prop:bspline}), the hazard rate and the Prony extraction of the atoms (Sections~\ref{sec:ch18:law} and~\ref{sec:ch18:reach}), the Liouville operator and the diffusion for P\'olya's kernel (Proposition~\ref{prop:biane:liouville}, Theorem~\ref{thm:biane:clock}), the Stieltjes criterion (Theorem~\ref{mw:thm:stieltjes}), the incomplete gamma form, the kink and Rouch\'e tracking (Proposition~\ref{mi:prop:incgamma}, Theorem~\ref{mi:thm:kink}, Corollary~\ref{mi:cor:rouche}). The results, with their status, are these.
\begin{enumerate}
\item Each term of the theta series for $\Phi$ is an explicit combination of two Gumbel kernels of fractional order, with Mellin transform $n^{-s}\xi(s)/(4\zeta(s))$ (Theorem~\ref{ev:thm:decomp}, proved). This gives the folded Gumbel kernels, P\'olya's two sided kernel and the Dyck chain, whose real zeros to the heights tabulated are computed (Table~\ref{ev:tab:approx}). The breakdown height $T_N\approx4(N+1)^2$ of the folded family is a heuristic balance of the kink against the size of $H$.
\item The zeros of $\Xi$ are all real and simple if and only if the reciprocal $\xi$ law is the hitting time of a string (Corollary~\ref{ev:cor:rh}, proved, with the hypotheses of simplicity and of the class of $\widehat H$ that the source left out). The tail gives $\gamma_1^2$ without any zero, and a Laplace inversion recovers $\gamma_1$ to eighteen digits (Proposition~\ref{ev:prop:tail}, proved under RH, simple zeros and $|\xi'(\rho)|\ge e^{-C|\gamma|}$ for the full series, the exponential excess needing only the two term tail; the digits are computed).
\item The Gumbel pieces have unconditional Hadamard products whose atoms are the poles of the archimedean factor; the atoms $\gamma_k^2$ come from the superposition, a sum and not a product (Section~\ref{ev:sec:bspline}, proved for the pieces, heuristic for the superposition).
\item P\'olya's kernel has an explicit string, the Liouville string of Biane, whose first eigenvalues are computed (Section~\ref{ev:sec:liouville}).
\item For a kernel with infinitely many zeros the zeros are all real if and only if the Hankel determinants of $\mu_j=\sum_k\widehat\gamma_k^{-2j}$ are all positive (Proposition~\ref{ev:prop:hankel}, proved). The first eleven pairs are positive for the exact kernel, the two term folded kernel and P\'olya's kernel (computed). Finite order positivity cannot see an off line pair above a height that grows with the order (heuristic).
\item Three routes to RH each stop at a known place, and the folded family reduces RH to a height rule for all $N$ (Section~\ref{ev:sec:path}, Proposition~\ref{ev:prop:reduction}, proved). The Rouch\'e inequality behind the rule is certified in ball arithmetic for $N=1,2,3,4$ at the heights of the table, given the count of zeros of $\Xi$ below the height (Section~\ref{ev:sec:work}); $N=5$ is computed.
\item The eigenvalue problem for the shifted Bessel kernels is a rational recurrence with real eigenvalues (proved), but the positivity of the extension and its reading as a string are not shown for the tabulated cases; the cotangent law for fractional shifts, the collisions of the de Bruijn Newman zeros and the masking of the first nonreal pair are in Sections~\ref{ev:sec:work} and~\ref{ev:sec:round2}, with the status stated there.
\item Devroye's alternating series gives a sign test for the reciprocal $\Xi$ law, the mean excess function of the hitting time has an expansion with a rate, and a detection rule for the first nonreal pair in the de Bruijn Newman family is stated and tested (Section~\ref{ev:sec:devroye}, computed unless stated).
\end{enumerate}
The remaining open problems close the chapter.
